\documentclass[10pt, letterpaper]{article}

% Packages:
\usepackage[
    ignoreheadfoot, % set margins without considering header and footer
    top=2 cm, % seperation between body and page edge from the top
    bottom=2 cm, % seperation between body and page edge from the bottom
    left=2 cm, % seperation between body and page edge from the left
    right=2 cm, % seperation between body and page edge from the right
    footskip=1.0 cm, % seperation between body and footer
    % showframe % for debugging 
]{geometry} % for adjusting page geometry
\usepackage{titlesec} % for customizing section titles
\usepackage{tabularx} % for making tables with fixed width columns
\usepackage{array} % tabularx requires this
\usepackage[dvipsnames]{xcolor} % for coloring text
\definecolor{primaryColor}{RGB}{0, 79, 144} % define primary color
\usepackage{enumitem} % for customizing lists
\usepackage{fontawesome5} % for using icons
\usepackage{amsmath} % for math
\usepackage[
    pdftitle={Michael Hanif Khan's CV},
    pdfauthor={Michael Khan},
    pdfcreator={LaTeX with RenderCV},
    colorlinks=true,
    urlcolor=primaryColor
]{hyperref} % for links, metadata and bookmarks
\usepackage[pscoord]{eso-pic} % for floating text on the page
\usepackage{calc} % for calculating lengths
\usepackage{bookmark} % for bookmarks
\usepackage{lastpage} % for getting the total number of pages
\usepackage{changepage} % for one column entries (adjustwidth environment)
\usepackage{paracol} % for two and three column entries
\usepackage{ifthen} % for conditional statements
\usepackage{needspace} % for avoiding page brake right after the section title
\usepackage{iftex} % check if engine is pdflatex, xetex or luatex

% Ensure that generate pdf is machine readable/ATS parsable:
\ifPDFTeX
    \input{glyphtounicode}
    \pdfgentounicode=1
    % \usepackage[T1]{fontenc} % this breaks sb2nov
    \usepackage[utf8]{inputenc}
    \usepackage{lmodern}
\fi



% Some settings:
\AtBeginEnvironment{adjustwidth}{\partopsep0pt} % remove space before adjustwidth environment
\pagestyle{empty} % no header or footer
\setcounter{secnumdepth}{0} % no section numbering
\setlength{\parindent}{0pt} % no indentation
\setlength{\topskip}{0pt} % no top skip
\setlength{\columnsep}{0cm} % set column seperation
\makeatletter
\let\ps@customFooterStyle\ps@plain % Copy the plain style to customFooterStyle
\patchcmd{\ps@customFooterStyle}{\thepage}{
    \color{gray}\textit{\small Michael Khan - Page \thepage{} of \pageref*{LastPage}}
}{}{} % replace number by desired string
\makeatother
\pagestyle{customFooterStyle}

\titleformat{\section}{\needspace{4\baselineskip}\bfseries\large}{}{0pt}{}[\vspace{1pt}\titlerule]

\titlespacing{\section}{
    % left space:
    -1pt
}{
    % top space:
    0.3 cm
}{
    % bottom space:
    0.2 cm
} % section title spacing

\renewcommand\labelitemi{$\circ$} % custom bullet points
\newenvironment{highlights}{
    \begin{itemize}[
        topsep=0.10 cm,
        parsep=0.10 cm,
        partopsep=0pt,
        itemsep=0pt,
        leftmargin=0.4 cm + 10pt
    ]
}{
    \end{itemize}
} % new environment for highlights

\newenvironment{highlightsforbulletentries}{
    \begin{itemize}[
        topsep=0.10 cm,
        parsep=0.10 cm,
        partopsep=0pt,
        itemsep=0pt,
        leftmargin=10pt
    ]
}{
    \end{itemize}
} % new environment for highlights for bullet entries


\newenvironment{onecolentry}{
    \begin{adjustwidth}{
        0.2 cm + 0.00001 cm
    }{
        0.2 cm + 0.00001 cm
    }
}{
    \end{adjustwidth}
} % new environment for one column entries

\newenvironment{twocolentry}[2][]{
    \onecolentry
    \def\secondColumn{#2}
    \setcolumnwidth{\fill, 4.5 cm}
    \begin{paracol}{2}
}{
    \switchcolumn \raggedleft \secondColumn
    \end{paracol}
    \endonecolentry
} % new environment for two column entries

\newenvironment{header}{
    \setlength{\topsep}{0pt}\par\kern\topsep\centering\linespread{1.5}
}{
    \par\kern\topsep
} % new environment for the header

\newcommand{\placelastupdatedtext}{% \placetextbox{<horizontal pos>}{<vertical pos>}{<stuff>}
  \AddToShipoutPictureFG*{% Add <stuff> to current page foreground
    \put(
        \LenToUnit{\paperwidth-2 cm-0.2 cm+0.05cm},
        \LenToUnit{\paperheight-1.0 cm}
    ){\vtop{{\null}\makebox[0pt][c]{
    }}}%
  }%
}%

% save the original href command in a new command:
\let\hrefWithoutArrow\href

% new command for external links:
\renewcommand{\href}[2]{\hrefWithoutArrow{#1}{\ifthenelse{\equal{#2}{}}{ }{#2 }\raisebox{.15ex}{\footnotesize \faExternalLink*}}}


\begin{document}
\newcommand{\AND}{\unskip
	\cleaders\copy\ANDbox\hskip\wd\ANDbox
	\ignorespaces
}
\newsavebox\ANDbox
\sbox\ANDbox{}

\placelastupdatedtext
\begin{header}
	\textbf{\fontsize{24 pt}{24 pt}\selectfont Michael Hanif Khan}

	\vspace{0.3 cm}

	\normalsize
	\mbox{{\color{black}{\footnotesize\faEnvelope[regular]}\hspace*{0.13cm}mkhan398 at uic dot edu}}%
	\kern 0.25 cm%
	\AND%
	\kern 0.25 cm%
	\mbox{\hrefWithoutArrow{https://xmasscan.github.io}{\color{black}{\footnotesize\faLink}\hspace*{0.13cm}xmasscan.github.io}}%
	\kern 0.25 cm%
	\AND%
	\kern 0.25 cm%
	\mbox{\hrefWithoutArrow{https://linkedin.com/in/mkhan398}{\color{black}{\footnotesize\faLinkedinIn}\hspace*{0.13cm}mkhan398}}%
	\kern 0.25 cm%
	\AND%
	\kern 0.25 cm%
	\mbox{\hrefWithoutArrow{https://github.com/xmasscan}{\color{black}{\footnotesize\faGithub}\hspace*{0.13cm}xmasscan}}%
\end{header}

\vspace{0.3 cm - 0.3 cm}

\section{Education}

\begin{twocolentry}{
		\textit{August 2023 – Current}}
	\textbf{University of Chicago at Illinois (UIC)}

	\textit{Bachelor of Science (B.S.), Computer Science}
\end{twocolentry}

\vspace{0.10 cm}
\begin{onecolentry}
	\textbf{Coursework:} Data Structures, Machine Organization, Computational Theory, Software Design, Systems Programming, Algorithms, Principles of Concurrency
\end{onecolentry}

\section{Experience}

\textbf{Systems and Internet Security Lab - SISL @ UIC}
\begin{twocolentry}{
    \textit{Fall 2025 - Current}
    }
    \textbf{Research Assistant - EV Charging Protocol Fuzzer}
\end{twocolentry}
\begin{onecolentry}
    \begin{highlights}
      \item Implemented expanded protocol violation detection systems, decreasing time required for manual forensic analysis tenfold. 
      \item Deployed automatic fuzzing instrumentation, simultaneously improving testing in development and production environments.
      \item Orchestrated dynamically generated Replay and Payload Injection attacks with a fuzzer to evaluate a system's resistance to common cyberattacks. 
    \end{highlights}
\end{onecolentry}


\section{Research}

\textbf{MSCS Undergraduate Research Laboratory - MURL}

\begin{twocolentry}{
    \textit{Spring 2025}
    }
    \textbf{Rational Points on Ellpitic Curves}
\end{twocolentry}
\begin{onecolentry}
    \begin{highlights}
      \item Wrote a program to model basic arithmetical operations upon Ellpitic Curves in Python.
      \item Utilized MatPlotLib to model the complexity of Elliptic Curves and the computational power required to determine it.
    \end{highlights}
\end{onecolentry}


\section{Professional Association}
\textbf{Linux User Group}

\begin{twocolentry}{
    Fall 2025 - Current
    }

  President, CyberForce Team Captain
\end{twocolentry}

\begin{onecolentry}
	\textbf{Department of Energy's CyberForce}
	\begin{highlights}
    \item Coordinated an effort to harden a cloud-based network with Free, Open Source network analysis software, discovering 36 major structural vulnerabilties. 
    \item Solved CTF challenges under a tight time constraint with unfamiliar technologies, such as OpenCV, Universal Radio Hacker, and OpenPLC.
    \item Performed real time network and analysis during a simulated cyberattack to identify, prevent, and properly record vulnerabilties within a live industrial environment.
	\end{highlights}

  \textbf{LUG Events}
  \begin{itemize}
    \item Intro to Web Security: SQL Injections \& XSS 
    \item Linux Week
  \end{itemize}
\end{onecolentry}

\begin{twocolentry}{

		Fall 2023 - Current}
	\textbf{Association for Computing Machinery (ACM)}
\end{twocolentry}
\begin{onecolentry}
  \begin{highlights}
    \item Founded SIG Hackathon, a student group centered around training students to compete in hackathons.
    \item HackMIT 2024 Participant
  \end{highlights}
\end{onecolentry}

\section{Technical Skills}

\begin{onecolentry}
	\textbf{Languages:} C, C++, Java, Python, F\#, Rust, Go, SQL
\end{onecolentry}

\vspace{0.2 cm}

\begin{onecolentry}
	\textbf{Technologies:} Linux, GitHub, Git, Docker, Proxmox, GDB, Valgrind, Makefile, AWS EC2
\end{onecolentry}

\vspace{0.2 cm}

\end{document}
